\section{Fixed-height measurements and the noncentral zeros}\label{sec:fixed}

\subsection{A signed three-height arithmetic estimate}

At the parameters \eqref{eq:band}, consider
\begin{equation}\label{eq:threeheight}
 W_T(a)=2\log|H_T(a+6\ii)|
              -\log|H_T(a+4\ii)|-\log|H_T(a+8\ii)|.
\end{equation}
Normalize by \(V=H_T/C_-\), and write
\(\nu=a/2+\pi T/8\).

\begin{lemma}[Complete outer comparison]\label{lem:outer}
For \(y=4,6,8\), the inherited field input implies, after a finite
uniform onset,
\begin{equation}\label{eq:outer}
 \left|V(a+\ii y)-
 \zeta\left(\frac{1+y}{2}+\frac c4-\ii\nu\right)\right|
 \le\frac{C_F^*}{L}.
\end{equation}
The constant \(C_F^*\) is finite and independent of \(a,c\).
\end{lemma}
\begin{proof}
Divide \eqref{eq:FN} by \(C_-\). The dominant exponential arm is a
Dirichlet sum whose argument is
\((1-\ii z)/2+T L(z)/4-\ii\pi T/8\), with the additional quadratic
weight \(\exp(T(\log n)^2/4)\).
The calculation in GNP-RS E01 \cite{gnpe01}, extended over the finite
set \(\{4,6,8\}\), has the following complete ledger:
\[
\begin{array}{ll}
\text{reflected population}&O(a^{(1-y)/4}),\\
\text{retained correction}&O(a^{-(y+1)/4}),\\
\text{complete normalized field error}&O(a^{-3/4}),\\
\text{quadratic-weight removal}&\le T D_2/4,\\
\text{ordinary Dirichlet tail}&\le(2/3)N^{-3/2},\\
\text{argument replacement}&\le D_1T|L(a+\ii y)-L|/4,
\end{array}
\]
where \(D_j=\sum_{n\ge2}n^{-5/2}(\log n)^j<\infty\).
For the weight removal, \(T\log n/4\le c/8\), and the elementary
bound \(|\ee^v-1|\le v\ee^v\) leaves a summable majorant.
Every other term is \(O(1/L)\) or smaller, uniformly at the three
heights. This yields \eqref{eq:outer}; the population is summed before
its error is bounded.
\end{proof}

Since
\[
 |\zeta(s)-1|\le 2^{-5/2}+\int_2^\infty x^{-5/2}\dd x
               =\frac7{12\sqrt2}<\frac12
 \quad(\Rea s\ge5/2),
\]
the gate \(L>4C_F^*\) supplies \(|V|>1/4\). Integration of the
logarithmic derivative along the straight segment from \(\zeta(s)\) to
\(V\) gives
\begin{equation}\label{eq:outerlog}
 \bigl|\log|V|-\log|\zeta(s)|\bigr|\le4C_F^*/L.
\end{equation}
This is an off-axis floor at these observations only.

Let \(g(z)=\log C_-(T,z)\), with the local analytic logarithm. Then
\begin{align}
 g'(z)&=\frac7{4z}-\frac\pi8+\frac{TL(z)}{8z}
              -\frac{\ii L(z)}4-\frac{\ii\pi T}{16z},\label{eq:gprime}\\
 g''(z)&=-\frac7{4z^2}+\frac{T(1-L(z))}{8z^2}
              -\frac\ii{4z}+\frac{\ii\pi T}{16z^2}.
\end{align}
For \(u(y)=\log|C_-(T,a+\ii y)|\), \(u''=-\Rea g''\).
For large \(a\), \(T|L(a+\ii y)|\le2\), and
\(|\Rea(-\ii/(a+\ii y))|=y/(a^2+y^2)\). Consequently
\[
 |u''(y)|\le5/a^2\quad(4\le y\le8),\qquad
 |D_C|\le20/a^2,\quad D_C=2u(6)-u(4)-u(8).
\]

\begin{proposition}[Three-height bound]\label{prop:W}
Let \(q_n=\Lambda_{\vM}(n)/\log n\), where \(\Lambda_{\vM}\) is the
von Mangoldt function. Then
\begin{equation}\label{eq:Wseries}
 W_T(a)=-\sum_{n\ge2}q_n n^{-5/2-c/4}(1-n^{-1})^2
                    \cos(\nu\log n)+D_C+e,\qquad
 |e|\le16C_F^*/L.
\end{equation}
In particular,
\begin{equation}\label{eq:Wbound}
 |W_T(a)|<\frac16+\frac{16C_F^*}{L}+\frac{20}{a^2}.
\end{equation}
Under the eventual gates
\begin{equation}\label{eq:Wgates}
 16C_F^*/L+20/a^2\le1/72,\qquad a^2\ge288,
\end{equation}
one has \(|W_T(a)|<13/72\).
\end{proposition}
\begin{proof}
The absolutely convergent Euler logarithm of ordinary zeta in
\(\Rea s>1\) \cite{dlmfzeta} is
\(\log\zeta(s)=\sum_{n\ge2}q_n n^{-s}\).
Combine its three values before taking absolute values. Their
coefficient is \(2n^{-1}-1-n^{-2}=-(1-n^{-1})^2\), with the same
phase \(\nu\log n\) throughout. The four total units of logarithmic
error in \eqref{eq:outerlog} give \eqref{eq:Wseries}.
The function \(f(x)=x^{-5/2}(1-x^{-1})^2\) decreases for \(x\ge2\);
since \(0\le q_n\le1\),
\[
 \sum_{n\ge2}q_n n^{-5/2-c/4}(1-n^{-1})^2
 \le f(2)+\int_2^\infty f(x)\dd x
 =\frac{389}{1680\sqrt2}<\frac16.
\]
This proves the claimed bound. A truncation at an integer \(M\ge2\)
would have tail at most
\(M^{-3/2-c/4}/(3/2+c/4)\); no truncation or sampled value is needed.
\end{proof}

\subsection{Why the scalar bound requires a negative background}

At an all-real slice the same three-height combination of the product is
\begin{equation}\label{eq:K}
 W_T(a)=\sum_r K(a-r),\qquad
 K(d)=\frac12\log\frac{(d^2+36)^2}{(d^2+16)(d^2+64)}.
\end{equation}
The sum converges absolutely because \(K(d)=-4d^{-2}+O(d^{-4})\).
At an event of multiplicity \(m\ge2\),
\begin{equation}
 W_T(a)=m[K(0)+K(2a)]+\mathcal R_{\mathrm{bg}},\qquad
 \mathcal R_{\mathrm{bg}}=\sum_{r\ne\pm a}K(a-r).
 \label{eq:Wbackground}
\end{equation}
Now \(K(0)=\log(9/8)>1/9\). For \(d^2\ge64\), direct comparison of
the numerator and denominator gives \(|K(d)|\le8/d^2\).
Thus \eqref{eq:Wgates} implies
\[
 m[K(0)+K(2a)]>2(1/9-1/144)=15/72.
\]
Combining with Proposition~\ref{prop:W} gives the quantitative consequence
\begin{equation}\label{eq:bgnegative}
 \boxed{\mathcal R_{\mathrm{bg}}<-1/36.}
\end{equation}
It is a necessary signed background contribution under the event,
not a contradiction: \(K\) changes sign.

\begin{proposition}[Cancellation forces a sign-changing log kernel]
\label{prop:kernel}
Let \(b_j>0\) and \(\alpha_j\in\R\) be a finite collection, and put
\[
 K_\alpha(d)=\frac12\sum_j\alpha_j\log(d^2+b_j^2).
\]
If \(\sum_j\alpha_j=0\), this kernel is integrable and
\(\int_\R K_\alpha(d)\dd d=\pi\sum_j\alpha_j b_j\).
If also \(\sum_j\alpha_j b_j=0\) and \(K_\alpha(0)>0\), then
\(K_\alpha\) is negative somewhere.
\end{proposition}
\begin{proof}
The first cancellation removes the \(\log d^2\) term, leaving
\(O(d^{-2})\). For positive \(b,b_0\), differentiating in \(b\) gives
\[
 \int_\R\log\frac{d^2+b^2}{d^2+b_0^2}\dd d=2\pi(b-b_0).
\]
Form the linear combination to obtain the integral. A continuous,
nonnegative, integrable function that is positive at zero cannot have
zero integral.
\end{proof}

For \eqref{eq:K},
\[
 K(d)>0\ (|d|<\sqrt{34}),\qquad
 K(d)<0\ (|d|>\sqrt{34}),\qquad \int_\R K(d)\dd d=0.
\]
The proposition concerns finite fixed-height logarithmic combinations
that cancel the linear carrier. It does not rule out a signed arithmetic
estimate of the background, explicit deflation, or a different observable.
The genuine positive-source threshold control in
Section~\ref{sec:controls} satisfies the same scalar cap.

\subsection{A correlated matrix test}

Write \(h=-H_T'/H_T\), and let \(z_j=a+\ii b_j\), with
\((b_1,b_2,b_3)=(4,6,8)\). On an all-real slice,
\begin{equation}\label{eq:Pick}
 M_{ij}=\frac{h(z_i)-\overline{h(z_j)}}{z_i-\bar z_j}
       =\sum_r\frac1{(z_i-r)(\bar z_j-r)}.
\end{equation}
This is a Gram matrix. Set
\[
 u_j=\frac1{\ii b_j},\qquad v_j=\frac1{2a+\ii b_j},\qquad
 E=uu^*+vv^*.
\]
An event of multiplicity \(m\) necessarily satisfies
\begin{equation}\label{eq:Pickevent}
 M-mE\succeq0.
\end{equation}
This retains correlations among all three observations.

There is a direct link to the previous scalar background. Deflate the
event in the logarithmic derivative:
\[
 h_{\mathrm{bg}}(z)=h(z)+\frac m{z-a}+\frac m{z+a},\qquad
 p_{\mathrm{bg}}(y)=\Ima h_{\mathrm{bg}}(a+\ii y)\ge0.
\]
Integration of \(\partial_y\log|H_T(a+\ii y)|=\Ima h(a+\ii y)\)
gives
\[
 \mathcal R_{\mathrm{bg}}
   =\int_4^6 p_{\mathrm{bg}}(y)\dd y-\int_6^8p_{\mathrm{bg}}(y)\dd y.
\]
Positivity alone does not order these two integrals.

To evaluate \(M\) arithmetically, use disks
\(|z-z_j|\le1/2\) in the fixed provider rectangle. Set
\[
 Z(z)=\zeta\left(\frac{1-\ii z}{2}+\frac c4-\frac{\ii\pi T}{8}\right),
 \qquad \Delta=V-Z.
\]
The complete outer ledger above, now on \(7/2\le\Ima z\le17/2\),
has reflected cost \(O(a^{-5/8})\), correction \(O(a^{-9/8})\),
field error \(O(a^{-3/4})\), summable weight-removal cost \(O(1/L)\),
tail \(O(N^{-5/4})\), and argument cost \(O(1/(aL))\). Hence
\begin{equation}\label{eq:deltaderivative}
 \norm{\Delta}_{|z-z_j|\le1/2}\le C_\Delta/L,\qquad
 |\Delta'(z_j)|\le2C_\Delta/L.
\end{equation}
The second bound is Cauchy's estimate for a holomorphic difference.
Since \(\Rea s\ge9/4\),
\[
 |Z-1|\le \frac{13}{5\,2^{9/4}}<\frac34,\qquad
 |Z'|\le\frac12\sum_{n\ge2}n^{-9/4}\log n=:B_Z.
\]
The gate \(C_\Delta/L\le1/8\) gives \(|V|>1/8\), and therefore
\begin{equation}\label{eq:logdererror}
 \left|\frac{V'}V-\frac{Z'}Z\right|
 \le\frac{C_h}{L},\qquad C_h=(16+32B_Z)C_\Delta.
\end{equation}
All denominators are paid before the matrix is formed.

Put \(\sigma_j=(1+b_j)/2+c/4\) and
\(P_j=\sum_{n\ge2}\Lambda_{\vM}(n)n^{-\sigma_j}\ee^{\ii\nu\log n}\).
Then
\begin{align}
 h(z_j)&=-g'(z_j)-\frac\ii2 P_j+O(1/L),\\
 M_{ij}&=\frac{-g'(z_i)+\overline{g'(z_j)}}{\ii(b_i+b_j)}
       -\frac{P_i+\overline{P_j}}{2(b_i+b_j)}
       +O\!\left(\frac1{L(b_i+b_j)}\right).
 \label{eq:signedM}
\end{align}
The bounded prime-power sums and \eqref{eq:gprime} imply
\begin{equation}\label{eq:Mbulk}
 M=\frac L2 C+O_{\op}(1),\qquad C_{ij}=\frac1{b_i+b_j}.
\end{equation}
This is a statement about the exact canonical samples at every large
center in the band, whether or not it is an event. The leading matrix
is positive definite because
\[
 \alpha^*C\alpha=\int_0^\infty
             \left|\sum_j\alpha_j\ee^{-b_jt}\right|^2\dd t>0
 \quad(\alpha\ne0).
\]

\begin{proposition}[The fixed-height matrix admits fixed multiplicities]
\label{prop:matrixcap}
For large \(a\), the largest formal nonnegative event weight allowed by
\eqref{eq:Pickevent} is
\begin{equation}\label{eq:matrixcap}
 \mathfrak m(a,c)=
 \frac1{\lambda_{\max}(M^{-1/2}EM^{-1/2})}
 =\frac6{13}L+O(1).
\end{equation}
For every fixed integer \(m\ge1\), the exact matrix \(M-mE\) is positive
definite at every sufficiently large center in the band.
\end{proposition}
\begin{proof}
Let \(r_j=1/b_j\). For any distinct positive heights, the Cauchy matrix
satisfies
\begin{equation}\label{eq:Cauchyinverse}
 r^tC^{-1}r=2\sum_j b_j^{-1}.
\end{equation}
To prove it, set \(x=C^{-1}r\) and \(f(z)=\sum_jx_j/(z+b_j)\).
Since \(f(b_i)=1/b_i\),
\[
 1-zf(z)=(-1)^q\prod_{j=1}^q\frac{z-b_j}{z+b_j}.
\]
Differentiate at zero. The left derivative is \(-f(0)\); the right
derivative is \(-2\sum b_j^{-1}\), which proves the identity.

Equation~\eqref{eq:Mbulk} gives
\(M^{-1}=2C^{-1}/L+O_{\op}(L^{-2})\).
The \(u\) term has its single nonzero eigenvalue
\(u^*M^{-1}u=4\sum b_j^{-1}/L+O(L^{-2})\).
The \(v\) term has norm \(O(a^{-2}/L)\). Since
\(\sum_{j=1}^3b_j^{-1}=13/24\), inversion proves
\eqref{eq:matrixcap}. Finally
\[
 \lambda_{\min}(M-mE)
 \ge \tfrac L2\lambda_{\min}(C)-O(1)-m\norm{E}>0
\]
for fixed \(m\) and large \(L\).
\end{proof}

This does not construct a zero configuration realizing the formal event
weight. It shows that this necessary matrix condition cannot reject a
fixed multiplicity in the tested regime. Its scale has a simple
interpretation:
\[
 \int_\R\frac{\dd r}{(a+\ii b_i-r)(a-\ii b_j-r)}
       =\frac{2\pi}{b_i+b_j}.
\]
A continuous density \(L/(4\pi)\) produces exactly the bulk \(LC/2\).
The fixed observations retain a growing background while the proposed
event has fixed weight. Substituting \(b_j=\beta_j/L\) into
\eqref{eq:Mbulk} is not justified by its fixed-geometry error estimate;
the microscopic problem requires a new acquisition argument.
